\documentclass[10pt,a4paper]{amsart}
\usepackage[margin=19mm]{geometry}
\usepackage{amsmath,amssymb,mathtools}
\usepackage[scr=rsfs,cal=euler]{mathalfa}
\usepackage{xfrac}
\usepackage[unicode,hidelinks,bookmarks=false]{hyperref}
\newcommand{\ie}{i.e.\ }
\newcommand{\cf}{cf.\ }
\newcommand{\resp}{respectively\ }
\newcommand{\Subsection}{\subsection}
\providecommand{\nobreakdash}{\nobreak\hbox{-}}
\setlength{\emergencystretch}{2em}
\multlinegap0pt
\usepackage{bm}
\usepackage{xcolor}
\usepackage{mathtools}
\usepackage[shortlabels]{enumitem}
\usepackage{algorithm}\usepackage{algpseudocode}
\usepackage{tikz}
\newtheorem{assumption}{Assumption}
\newtheorem{lemma}{Lemma}
\usepackage{cleveref}
\crefname{assumption}{Assumption}{Assumptions}
\crefname{lemma}{Lemma}{Lemmas}
\newcommand{\E}{\mathbb{E}}
\newcommand{\abs}[1]{\left|#1\right|}
\newcommand{\blipadv}{\tau}
\newcommand{\norm}[1]{\lVert#1\rVert}
\renewcommand{\qed}{\unskip\nobreak\hfill\ensuremath{\square}}
\title{Decision-Centered Abstractions via Orthogonal Estimation of Difference-of-Q Functions}
\author{Defu Cao, Angela Zhou}
\date{Source: arXiv:2406.08697v4 (CC BY 4.0)}
\begin{document}
\maketitle

\noindent\emph{Source and licence.} Decision-Centered Abstractions via Orthogonal Estimation of Difference-of-Q Functions. Version arXiv:2406.08697v4 (CC BY 4.0), distributed by arXiv under the Creative Commons Attribution 4.0 International licence, http://creativecommons.org/licenses/by/4.0/, as recorded in the arXiv licence metadata for this exact version and shown on the abstract page https://arxiv.org/abs/2406.08697v4. Full licence notice: \texttt{licenses/arxiv-2406.08697-CC-BY-4.0.txt}.\par\smallskip

\noindent\textit{Editorial scope.} Counted unit: the article's lemma lemma-adverror-to-value (source0.tex lines 2338--2359). The article's complete original proof of it is delivered with it (source0.tex lines 3874--3945, one \texttt{proof} environment of 72 lines, which the article places away from the statement). No mathematics is written, repaired or re-derived: the statement and the proof are the article's own lines, in the article's own notation, and no retained line is substituted or re-flowed.  Retained uncounted as the source-local context the proof needs: lines 2292--2317.  Nothing was omitted from any retained span, so the package carries no omission record.  No retained line contains a citation, so the wrapper carries no bibliography and no bibliography entry.  No retained line contains a figure environment, an image-inclusion command or drawing code, so no image is delivered and the package declares no asset dependency.  Editorial formatting only, in addition: an AMS \texttt{amsart} preamble that restates the article's own declarations and macros used by the retained lines, namely the environments \texttt{assumption}, \texttt{lemma} on separate counters, together with the standard packages those lines need.  The retained environments print in this document as Assumption 1, Lemma 1 in order of appearance, whereas the article prints the same items with the numbering of its own full document.  The complete native source of the article as distributed in the arXiv e-print for this exact version is bundled unchanged as \texttt{sources/160733/source0.tex}.
\begin{assumption}[Margin on observational distribution]\label{asn-margin-observational-pib}
% There exist constants $0<\eta \leq M$ and $0<\alpha<\infty$ such that
% $$
% P_X(|\blipadv(X)| \leq t) \leq(t / \eta)^\alpha, \quad \forall 0 \leq t \leq \eta,
% $$
% where $M<\infty$ is the constant as defined in Assumption 2.1 (BO).
Let $Q^{*}_t(s, \pi^*)$ be the optimal $Q$ function at the optimal action, and $a^{\prime}$ the second-best option, $a^{\prime} \in \mathcal{A}\setminus \arg \max _a Q^{*}_t(s, a).$
Assume there exist constants $\alpha\geq0$ (the margin exponent) and
$\delta_0>0$ such that
$$
\begin{aligned}
\textstyle
% & \lambda\left\{s \in \mathbb{S}: \max _a Q^{\text {opt }}(s, a)-\max _{a^{\prime} \in \mathcal{A}\setminus \arg \max _a Q^{o p t}(x, a)} Q^{\text {opt }}\left(x, a^{\prime}\right) \leq \epsilon\right\}
P_{ 
% d^{\pi^b_t}
}
\left( Q^{*}_t(S_t, \pi^*)
%\max _a Q^{*}_t(s, a)
-Q^{*}_t\left(S_t, a^{\prime}\right) \leq \epsilon\right)
% =O\left(
\leq (\epsilon/\delta_0)^\alpha
% \right)
, \forall t \in 1, \dots, T
\end{aligned}
$$
\end{assumption}

\begin{lemma}[Contrast error to decision regret under a margin condition]
\label{lemma-adverror-to-value}
Fix a stage $t$ and binary actions. Suppose \Cref{asn-margin-observational-pib} holds with
exponent $\alpha\geq0$. For any event $\mathcal E$ on which
$\norm{\hat\blipadv_t-\blipadv_t}_2\leq\epsilon_t$, the following bounds hold
on the same event:
\begin{equation*}
\E\!\left[
Q_t^*(S_t,\pi_t^*(S_t))-Q_t^*(S_t,\hat\pi_t(S_t))
\right]
\lesssim
\epsilon_t^{\frac{2+2\alpha}{2+\alpha}}.
\end{equation*}
\begin{equation*}
\norm{
Q_t^*(S_t,\pi_t^*(S_t))-Q_t^*(S_t,\hat\pi_t(S_t))
}_2
\leq\epsilon_t.
\end{equation*}
% The expectation and norm are over the observational stage-$t$ state
% distribution, conditional on the fitted policy when appropriate.
\end{lemma}

\begin{proof}{Proof of \Cref{lemma-adverror-to-value}}
For binary actions, an action mistake implies
\[
Q_t^*(S_t,\pi_t^*(S_t))-Q_t^*(S_t,\hat\pi_t(S_t))
=\abs{\blipadv_t(S_t)}\mathbb I[\hat\pi_t(S_t)\ne\pi_t^*(S_t)]
\leq\abs{\hat\blipadv_t(S_t)-\blipadv_t(S_t)}.
\]

Suppose the event $\mathcal E$ supplies the integrated $L_2$ bound. For
$\epsilon>0$, keep
the mistake-region set
$\mathcal S_\epsilon=
\{s:\max_a Q_t^*(s,a)-Q_t^*(s,\hat\pi_t(s))\le\epsilon\}$.
The pointwise mistake bound and the disjoint partition give
\begin{align}
&\norm{
Q_t^*(S_t,\pi_t^*(S_t))-Q_t^*(S_t,\hat\pi_t(S_t))
}_2^2\notag\\
&\quad=
\norm{
\{Q_t^*(S_t,\pi_t^*(S_t))-Q_t^*(S_t,\hat\pi_t(S_t))\}
\mathbb I[S_t\in\mathcal S_\epsilon]
}_2^2\notag+
\norm{
\{Q_t^*(S_t,\pi_t^*(S_t))-Q_t^*(S_t,\hat\pi_t(S_t))\}
\mathbb I[S_t\notin\mathcal S_\epsilon]
}_2^2\notag\\
&\quad
\le
\norm{\hat\blipadv_t-\blipadv_t}_2^2.
\label{eqn-mistakeregion-ltwo}
\end{align}
For expected regret, the same partition yields
\begin{align*}
&\E[
Q_t^*(S_t,\pi_t^*(S_t))-Q_t^*(S_t,\hat\pi_t(S_t))
]\\
&\quad\leq
\epsilon
P\!\left(
0<\max_a Q_t^*(S_t,a)-Q_t^*(S_t,\hat\pi_t(S_t))
\le\epsilon
\right)
+
\E[
\abs{\hat\blipadv_t(S_t)-\blipadv_t(S_t)}
\mathbb I[\abs{\hat\blipadv_t(S_t)-\blipadv_t(S_t)}
>\epsilon]
]\\
&\quad\lesssim
\epsilon^{1+\alpha}
+\epsilon^{-1}
\norm{\hat\blipadv_t-\blipadv_t}_2^2.
\end{align*}
Choosing
$\epsilon=\norm{\hat\blipadv_t-\blipadv_t}_2^{2/(2+\alpha)}$
gives
\begin{align*}
&\E[
Q_t^*(S_t,\pi_t^*(S_t))-Q_t^*(S_t,\hat\pi_t(S_t))
] \lesssim
\norm{\hat\blipadv_t-\blipadv_t}_2^
{\frac{2+2\alpha}{2+\alpha}},\\
&\norm{
Q_t^*(S_t,\pi_t^*(S_t))-Q_t^*(S_t,\hat\pi_t(S_t))
}_2\leq
\norm{\hat\blipadv_t-\blipadv_t}_2.
\end{align*}
Substituting the eventwise bound
$\norm{\hat\blipadv_t-\blipadv_t}_2\leq\epsilon_t$
shows that both conclusions hold on the same event. \qed
\end{proof}

\end{document}
